\chapter{Fazit}

Das Projektmodul bot den Studierenden die Möglichkeit, theoretische Inhalte aus dem Studium in einem praxisnahen Umfeld anzuwenden und durch die Bearbeitung realer Aufgabenstellungen zu vertiefen. Die Zusammenarbeit mit dem Oldenburger Computer-Museum stellte dabei sicher, dass die entwickelten Projekte nicht nur einen akademischen Hintergrund hatten, sondern einen konkreten Anwendungsbezug besaßen.

Trotz unterschiedlicher Themen und Herangehensweisen verfolgten alle Projektgruppen das gemeinsame Ziel, historische Computertechnik auf moderne und anschauliche Weise zu vermitteln. Dabei entstanden vielfältige Lösungen, die von der Konzeption über die technische Umsetzung bis hin zur Dokumentation eigenständig entwickelt wurden.

Neben den fachlichen Ergebnissen förderte das Projektmodul insbesondere die Arbeit im Team, die eigenverantwortliche Organisation sowie den regelmäßigen Austausch innerhalb der Gruppen und mit den Projektbetreuern. Gleichzeitig bot das Modul den Freiraum, eigene Ideen einzubringen, verschiedene Lösungsansätze auszuprobieren und Erfahrungen im Umgang mit realen Projektabläufen zu sammeln.

Die in dieser Dokumentation vorgestellten Projekte verdeutlichen, wie durch die Verbindung von technischem Wissen, Kreativität und Teamarbeit innovative Konzepte entstehen können, die einen Mehrwert für das Oldenburger Computer-Museum schaffen und gleichzeitig einen praxisnahen Einblick in die Projektarbeit vermitteln.
